% from ega4-9-fr.tex
\oldpage[IV-3]{75}

Cela étant, il résulte de (\textbf{II}, 4.4.4) que dire que $\mathcal L_s$
n'est pas très ample relativement à $f_s$ signifie que~: \emph{ou bien}
l'homomorphisme canonique
$(f_s)^*((f_s)_*(\mathcal L_s))\to\mathcal L_s$ n'est pas surjectif~; \emph{ou
bien} l'homomorphisme précédent est surjectif et le morphisme canonique
$r:X_s\to\mathbf P((f_s)_*(\mathcal L_s))$ n'est pas une immersion. Compte
tenu de l'isomorphisme (9.6.3.1), ces conditions s'écrivent respectivement sous
la forme~: $1^\circ$ l'homomorphisme canonique
$(f^*(f_*(\mathcal L)))_s\to\mathcal L_s$ n'est pas surjectif~; $2^\circ$
l'homomorphisme précédent est surjectif et le morphisme canonique
$X_s\to\mathbf P((f_*(\mathcal L))_s)$ n'est pas une immersion.

Supposons d'abord que l'homomorphisme canonique
$(f^*(f_*(\mathcal L)))_\eta\to\mathcal L_\eta$ ne soit pas surjectif. Comme
$f_*(\mathcal L)$ est cohérent, il en est de même de $f^*(f_*(\mathcal L))$,
et alors il résulte de (9.4.5) que pour tout $s$ dans un voisinage de $\eta$,
l'homomorphisme $(f^*(f_*(\mathcal L)))_s\to\mathcal L_s$ n'est pas
surjectif, ce qui prouve dans ce cas que $S-E'$ est un voisinage de $\eta$.

Supposons en second lieu que l'homomorphisme canonique
$(f^*(f_*(\mathcal L)))_\eta\to\mathcal L_\eta$ soit surjectif mais que le
morphisme $X_\eta\to\mathbf P((f_*(\mathcal L))_\eta)$ ne soit pas une
immersion. Alors le même raisonnement que ci-dessus montre d'abord que pour
tout $s$ assez voisin de $\eta$, l'homomorphisme
$(f^*(f_*(\mathcal L)))_s\to\mathcal L_s$ est surjectif~; d'autre part, en
vertu de (9.6.1, (viii)), pour $s$ assez voisin de $\eta$, le morphisme
$X_s\to\mathbf P((f_*(\mathcal L))_s)$ n'est pas une immersion, ce qui achève
la démonstration dans le cas de $E'$.

II) \emph{Cas de $E$.} --- Considérons d'abord le cas particulier où $Y=S$~:

\begin{corollary}[9.6.4]
\label{IV.9.6.4-fr}
Soit $f:X\to S$ un morphisme propre de présentation finie, $\mathcal L$ un
$\mathcal O_X$-Module inversible. Alors, l'ensemble $E$ des $s\in S$ tels que
$\mathcal L_s$ soit ample (relativement à $f_s$) est ouvert dans $S$, et
$\mathcal L|f^{-1}(E)$ est ample relativement à la restriction
$f^{-1}(E)\to E$ de $f$.
\end{corollary}

Comme la condition (9.2.1, (i)) est vérifiée par la propriété $\mathbf P$ définie plus
haut, le résultat de (9.2.2, (iv)) et le raisonnement de (9.2.3) montrent qu'on
peut se borner au cas où $S$ est \emph{noethérien}~; mais alors le résultat
découle de (\textbf{III}, 4.7.1) et de la stabilité de l'amplitude par
changement de base (\textbf{II}, 4.6.13).

\begin{corollary}[9.6.5]
\label{IV.9.6.5-fr}
Sous les hypothèses de (9.6.4), pour que $\mathcal L$ soit ample relativement
à $f$, il faut et il suffit que, pour tout $s\in S$, $\mathcal L_s$ soit ample
relativement à $f_s$.
\end{corollary}

\phantomsection\label{IV.9.6.6-fr}%
\textbf{(9.6.6) Fin de la démonstration de (9.6.3).} --- Revenons au cas
général, $S$ étant noethérien, intègre et de point générique
$\eta\in S-E$. Comme $f$ est propre, il résulte de (9.6.4) que l'ensemble $V$
des $y\in Y$ tels que $\mathcal L_y$ soit un $\mathcal O_{X_y}$-Module ample,
relativement au morphisme $f_y:X_y\to\operatorname{Spec}(k(y))$ est ouvert,
donc $F=Y-V$ est fermé dans $Y$. Cela étant, comme $f_s$ est propre et que,
pour tout $y\in Y$ au-dessus de $s\in S$, on a
$\mathcal L_y=(\mathcal L_s)_y$, il résulte de (9.6.5) que pour que $s$
appartienne à $S-E$, il faut et il suffit que l'on ait $F_s\neq\varnothing$.
Mais comme $F$ est fermé dans $Y$ et que $F_\eta\neq\varnothing$, il résulte
de (9.5.1) que l'ensemble des $s\in S$ tels que $F_s\neq\varnothing$ est un
voisinage de $\eta$ dans $S$. C.Q.F.D.

\begin{remarks}[9.6.7]
\label{IV.9.6.7-fr}
(i) Pour chacune des propriétés $\mathbf P$ considérées dans (9.6.1) la proposition
(9.3.3) est applicable, et ces propriétés (pour les morphismes $f_s$) sont donc
«~stables~» par passage d'une limite projective d'un système projectif
essentiellement affine (8.13.4) de préschémas $X_\lambda$ à l'un convenable
d'entre eux.

(ii) Soit $Z$ une partie \emph{localement constructible} de $X$ telle que, pour
tout $s\in S$, $Z\cap X_s$ soit \emph{ouvert} dans $X_s$, et désignons par
$Z_s$ le sous-préschéma de $X_s$ induit sur l'ouvert $Z\cap X_s$.

